%% kapa: exam=1 round=34 year=2023 period=1 form=A subject=economics subject_ko=경제학원론 seq=19 number=59
%% source: Q-Net 감정평가사 기출문제 (cst00309) artlSeq=5211789 file=raw/1008/5211789/2023년도 제34회 감정평가사 1차시험 1교시.pdf
%% conversion: kordoc 4.24.51 -> tools/gen1.py (auto); review=proofread-sonnet; flags=none
\begin{question}{59}
정보의 비대칭성에 관한 설명으로 옳은 것은?
\begin{choices}
  \item 도덕적 해이(moral hazard)는 감춰진 속성(hidden characteristics)과 관련된다.
  \item 직업감독제도는 역선택(adverse selection) 방지를 위한 효율적인 수단이다.
  \item 자동차보험에서 기초공제(initial deduction)제도를 두는 이유는 역선택 방지를 위함이다.
  \item 상품시장에서 역선택 방지를 위해 품질보증제도를 도입한다.
  \item 노동시장에서 교육수준을 선별의 수단으로 삼는 이유는 도덕적 해이를 방지하기 위함이다.
\end{choices}
\end{question}
\answer{4}
